\documentclass[]{article}

\usepackage{amsmath}
\usepackage{multirow}
\usepackage{natbib}
\usepackage{graphicx} 


\begin{document}

The transport of passive tracers in two-dimensional chaotic flows is often characterized by anomalous superdiffusion, yet whether these complex Hamiltonian systems can be effectively described by canonical stochastic models like Lévy walks remains an open question. We address this by directly comparing numerical simulations of tracer trajectories in point-vortex systems of varying chaoticity, controlled by the number of vortices $N$, with a benchmark dataset of Lévy walks. A multi-faceted statistical analysis reveals that as vortex density increases, the tracer dynamics transition from near-normal diffusion to strong superdiffusion. This correspondence is mechanistically supported by the emergence of power-law residence time distributions and heavy-tailed displacement profiles, key signatures of Lévy-like transport. Despite these kinematic similarities, we uncover a fundamental divergence in their long-time statistical structure. We demonstrate that the vortex system becomes progressively more ergodic as superdiffusion strengthens with increasing $N$, a trend that is diametrically opposed to the increasing non-ergodicity of superdiffusive Lévy walks. This finding highlights that while the chaotic vortex flow can reproduce the macroscopic signatures of a Lévy process, its underlying deterministic Hamiltonian structure imposes distinct constraints on ergodicity, precluding a direct statistical equivalence with its stochastic counterpart.

\end{document}
                